\documentclass[tate,paper=a5,fontsize=10pt]{jlreq}
\usepackage{luatexja-ruby}
\begin{document}
\begin{center}{\Large 港の朝}\end{center}

\section*{一}
港町の朝は、いつも\ruby{汽笛}{きてき}の音で始まる。坂の上の下宿から海を見下ろすと、灰色の水面に白い船体がゆっくりと滑り込んでくるのが見えた。港の空には薄い雲が流れ、遠くの山の稜線がかすんで見えた。風向きが変わるたびに、潮の香りと油の匂いが交互に鼻をかすめた。通りの角では、自転車に乗った郵便配達夫がベルを鳴らして通り過ぎていった。僕はそのたびに立ち止まり、何かを思い出そうとしては、結局何も思い出せずに歩き出した。

「また今日も早いのね」と、下の階から\ruby{女将}{おかみ}の声がした。

「ええ、市場に寄ってから学校へ行きます」

僕は階段を下りながら答えた。女将は\ruby{割烹着}{かっぽうぎ}の袖をまくり、竈の火を見ていた。湯気の向こうで、味噌汁の匂いが立ちのぼっていた。

市場は人でごった返していた。魚屋の店先には\ruby{鰯}{いわし}や\ruby{鯵}{あじ}が銀色に光って並び、威勢のよい掛け声があちこちから飛んできた。僕はその間を縫うように歩き、いつもの乾物屋の前で足を止めた。港の空には薄い雲が流れ、遠くの山の稜線がかすんで見えた。風向きが変わるたびに、潮の香りと油の匂いが交互に鼻をかすめた。通りの角では、自転車に乗った郵便配達夫がベルを鳴らして通り過ぎていった。僕はそのたびに立ち止まり、何かを思い出そうとしては、結局何も思い出せずに歩き出した。

「坊主、今朝は何にする」

「\ruby{昆布}{こんぶ}を少しください。母に送るんです」

店の主人は黙ってうなずき、新聞紙で手際よく包んでくれた。代金を払うと、彼はおまけだと言って小さな飴玉を一つ掌に載せた。

\section*{二}
学校までの道は、海沿いの防波堤の上を通っている。潮風に吹かれながら歩くと、制服の襟がぱたぱたと鳴った。港の空には薄い雲が流れ、遠くの山の稜線がかすんで見えた。風向きが変わるたびに、潮の香りと油の匂いが交互に鼻をかすめた。通りの角では、自転車に乗った郵便配達夫がベルを鳴らして通り過ぎていった。僕はそのたびに立ち止まり、何かを思い出そうとしては、結局何も思い出せずに歩き出した。

防波堤の端に、見慣れない少女が立っていた。長い髪を風になびかせ、じっと沖の方を見つめている。僕が近づいても、彼女は振り向こうとしなかった。

「何を見ているの」と僕は尋ねた。

「\kenten{あの船}」と彼女は短く答え、水平線の一点を指さした。そこには小さな黒い影が浮かんでいるだけで、それが船なのかどうかも僕には分からなかった。

「あれに父が乗っているの。三年ぶりに帰ってくるはずなのよ」

彼女の横顔には、期待とも不安ともつかない表情が浮かんでいた。僕は何と言えばよいのか分からず、ただ黙って隣に立っていた。

やがて始業の鐘が遠くで鳴った。僕は慌てて走り出したが、振り返ると彼女はまだ同じ場所に立って、\ruby{微動}{びどう}だにせず海を見ていた。港の空には薄い雲が流れ、遠くの山の稜線がかすんで見えた。風向きが変わるたびに、潮の香りと油の匂いが交互に鼻をかすめた。通りの角では、自転車に乗った郵便配達夫がベルを鳴らして通り過ぎていった。僕はそのたびに立ち止まり、何かを思い出そうとしては、結局何も思い出せずに歩き出した。

\section*{三}
その日の授業は少しも頭に入らなかった。黒板に書かれた数式を写しながらも、防波堤の少女のことばかり考えていた。港の空には薄い雲が流れ、遠くの山の稜線がかすんで見えた。風向きが変わるたびに、潮の香りと油の匂いが交互に鼻をかすめた。通りの角では、自転車に乗った郵便配達夫がベルを鳴らして通り過ぎていった。僕はそのたびに立ち止まり、何かを思い出そうとしては、結局何も思い出せずに歩き出した。

放課後、僕はもう一度防波堤へ行ってみた。けれども彼女の姿はなく、ただ\ruby{鴎}{かもめ}が数羽、低く鳴きながら旋回しているだけだった。

沖を見ると、朝の黒い影は消えていた。代わりに一隻の\ruby{貨物船}{かもつせん}が、港の入り口に向かってゆっくりと進んでいる。港の空には薄い雲が流れ、遠くの山の稜線がかすんで見えた。風向きが変わるたびに、潮の香りと油の匂いが交互に鼻をかすめた。通りの角では、自転車に乗った郵便配達夫がベルを鳴らして通り過ぎていった。僕はそのたびに立ち止まり、何かを思い出そうとしては、結局何も思い出せずに歩き出した。

「あの船だろうか」と僕は\kenten{ひとりごと}を言った。

下宿に戻ると、女将が玄関先で待っていた。港の空には薄い雲が流れ、遠くの山の稜線がかすんで見えた。風向きが変わるたびに、潮の香りと油の匂いが交互に鼻をかすめた。通りの角では、自転車に乗った郵便配達夫がベルを鳴らして通り過ぎていった。僕はそのたびに立ち止まり、何かを思い出そうとしては、結局何も思い出せずに歩き出した。

「あんたに手紙が来ているよ。差出人の名前がないんだけどね」

封筒を受け取ると、中には一枚の便箋が入っていた。そこには見覚えのない筆跡で、ただ一行、明日の朝、防波堤でとだけ書かれていた。港の空には薄い雲が流れ、遠くの山の稜線がかすんで見えた。風向きが変わるたびに、潮の香りと油の匂いが交互に鼻をかすめた。通りの角では、自転車に乗った郵便配達夫がベルを鳴らして通り過ぎていった。僕はそのたびに立ち止まり、何かを思い出そうとしては、結局何も思い出せずに歩き出した。

\section*{四}
翌朝、僕はいつもより一時間早く家を出た。空はまだ薄暗く、街灯の明かりが濡れた石畳に映っていた。港の空には薄い雲が流れ、遠くの山の稜線がかすんで見えた。風向きが変わるたびに、潮の香りと油の匂いが交互に鼻をかすめた。通りの角では、自転車に乗った郵便配達夫がベルを鳴らして通り過ぎていった。僕はそのたびに立ち止まり、何かを思い出そうとしては、結局何も思い出せずに歩き出した。

防波堤に着くと、彼女はもうそこにいた。昨日と同じように沖を見つめていたが、今日は僕の足音に気づいて振り向いた。

「来てくれたのね」

「手紙をくれたのは君だったの」

彼女は小さくうなずいた。そして\ruby{外套}{がいとう}の内ポケットから、古びた写真を一枚取り出して僕に見せた。そこには若い船乗りと、幼い女の子が並んで写っていた。港の空には薄い雲が流れ、遠くの山の稜線がかすんで見えた。風向きが変わるたびに、潮の香りと油の匂いが交互に鼻をかすめた。通りの角では、自転車に乗った郵便配達夫がベルを鳴らして通り過ぎていった。僕はそのたびに立ち止まり、何かを思い出そうとしては、結局何も思い出せずに歩き出した。

「父よ。この写真の船が、今日の昼に入港するの。でも、一人で迎えに行くのが怖くて」

僕は写真をじっと見つめた。船乗りの胸には、\kenten{錨の形}をした小さなバッジが光っていた。港の空には薄い雲が流れ、遠くの山の稜線がかすんで見えた。風向きが変わるたびに、潮の香りと油の匂いが交互に鼻をかすめた。通りの角では、自転車に乗った郵便配達夫がベルを鳴らして通り過ぎていった。僕はそのたびに立ち止まり、何かを思い出そうとしては、結局何も思い出せずに歩き出した。

「一緒に行こう」と僕は言った。自分でも驚くほど、はっきりした声だった。

\section*{五}
昼過ぎ、港は出迎えの人々でにぎわっていた。\ruby{桟橋}{さんばし}の先に大きな船が横付けされ、\ruby{舷梯}{げんてい}から乗組員が次々と降りてくる。港の空には薄い雲が流れ、遠くの山の稜線がかすんで見えた。風向きが変わるたびに、潮の香りと油の匂いが交互に鼻をかすめた。通りの角では、自転車に乗った郵便配達夫がベルを鳴らして通り過ぎていった。僕はそのたびに立ち止まり、何かを思い出そうとしては、結局何も思い出せずに歩き出した。

彼女は僕の袖をつかみ、背伸びをして一人ひとりの顔を確かめていた。日に焼けた男たちが家族と抱き合い、笑い声や泣き声があちこちで上がった。

けれども、写真の船乗りの姿はなかなか現れなかった。最後の一人が降りきったとき、彼女の手から力が抜けるのが分かった。港の空には薄い雲が流れ、遠くの山の稜線がかすんで見えた。風向きが変わるたびに、潮の香りと油の匂いが交互に鼻をかすめた。通りの角では、自転車に乗った郵便配達夫がベルを鳴らして通り過ぎていった。僕はそのたびに立ち止まり、何かを思い出そうとしては、結局何も思い出せずに歩き出した。

「来なかった」と彼女はつぶやいた。

そのとき、船の\ruby{甲板}{かんぱん}から誰かが大きく手を振った。帽子を取ったその男の胸に、錨の形のバッジが光っているのが見えた。港の空には薄い雲が流れ、遠くの山の稜線がかすんで見えた。風向きが変わるたびに、潮の香りと油の匂いが交互に鼻をかすめた。通りの角では、自転車に乗った郵便配達夫がベルを鳴らして通り過ぎていった。僕はそのたびに立ち止まり、何かを思い出そうとしては、結局何も思い出せずに歩き出した。

「お父さん」

彼女は叫ぶと同時に走り出していた。僕はその背中を見送りながら、汽笛の音がいつもより長く、高く響いているように感じた。港の空には薄い雲が流れ、遠くの山の稜線がかすんで見えた。風向きが変わるたびに、潮の香りと油の匂いが交互に鼻をかすめた。通りの角では、自転車に乗った郵便配達夫がベルを鳴らして通り過ぎていった。僕はそのたびに立ち止まり、何かを思い出そうとしては、結局何も思い出せずに歩き出した。

\end{document}
